\section{A weighted threshold complex}

The Dirichlet series of $a$ is especially rigid.  For $\Re z>1$,
\begin{align}
 Q(z)&=\sum_{n\ge1}\frac{q(n)}{n^z}
 =6\zeta(z)-6+9\,2^{-z}-3\,4^{-z},\\
 A(z)&=\sum_{n\ge1}\frac{a(n)}{n^z}
 =\frac{Q(z)}{\zeta(z)}
 =6-\frac{3(1-2^{-z})(2-2^{-z})}{\zeta(z)}.
\end{align}
Let $d$ be defined by $6\delta_1-a=d$; then $\D(t)=\sum_{n\le t}d(n)/\sqrt n$.
Writing $z=s+1/2$ gives
\begin{equation}\label{eq:Euler-generalized}
 \frac16\sum_{n\ge1}\frac{d(n)}{n^{s+1/2}}
 =(1-2^{-s-1/2})^2(1-2^{-s-3/2})
  \prod_{p\,\mathrm{odd}}(1-p^{-s-1/2}).
\end{equation}

\begin{definition}[Generalized-prime labels]
There are two labels $2_1,2_2$ of cost $2$ and activity $2^{-1/2}$, one label
$2_3$ of cost $2$ and activity $2^{-3/2}$, and one label of cost $p$ and
activity $p^{-1/2}$ for every odd prime $p$.
For a finite labelled subset $S$, let $P(S)$ and $r(S)$ be the products of its
costs and activities.
\end{definition}

\begin{theorem}[Weighted Euler characteristic]\label{thm:Eulerchar}
For every $t\ge1$,
\begin{align}
 \frac{\D(t)}6&=\sum_{S:\,P(S)\le t}(-1)^{|S|}r(S),\label{eq:Z}\\
 \frac{C(t)}6&=\sum_{\varnothing\ne S:\,P(S)\le t}
                     (-1)^{|S|+1}r(S).\label{eq:F}
\end{align}
If each available label is independently retained with probability equal to
its activity and $K_t[R]$ is the induced threshold complex, then
\begin{equation}\label{eq:expectedchi}
 \boxed{\frac{C(t)}6=\mathbb E\,\chi(K_t[R]).}
\end{equation}
\end{theorem}
\begin{proof}
For a fixed $t$, only finitely many labels can occur.  Expanding the finite
product in \eqref{eq:Euler-generalized} gives \eqref{eq:Z} coefficient by
coefficient; \eqref{eq:F} follows from $C=6-\D$.  Finally,
\[
 \mathbb E\chi(K_t[R])
 =\sum_{\varnothing\ne S:\,P(S)\le t}(-1)^{|S|+1}\Pr(S\subseteq R),
\]
and independence gives $\Pr(S\subseteq R)=r(S)$.
\end{proof}

This reformulation is not a positivity proof: a threshold complex can have
negative Euler characteristic.  It does, however, identify the exact
orientation information absent from an unsigned energy estimate.

\section{Exact windowed Bellman recurrence}

For a finite label set $V$, set
\[
 Z_V(t)=\sum_{S\subseteq V,\,P(S)\le t}(-1)^{|S|}r(S),
 \quad Z_V(t)=0\ (t<1),
\]
\[
 F_V(t)=\1_{t\ge1}-Z_V(t),
 \qquad G_V(X)=\int_{X/4}^{X}F_V(t)\frac{dt}{t}.
\]

\begin{theorem}[One-label Bellman update]\label{thm:Bellman}
If a new label has cost $p$ and activity $r$, then
\begin{equation}\label{eq:Zrec}
 Z_{V\cup\{p\}}(t)=Z_V(t)-rZ_V(t/p),
\end{equation}
and
\begin{equation}\label{eq:Grec}
 \boxed{G_{V\cup\{p\}}(X)
 =G_V(X)+r\bigl(H_X(p)-G_V(X/p)\bigr)}.
\end{equation}
For the complete generalized-prime set, $G(X)=\A_X/6$.
\end{theorem}
\begin{proof}
Split subsets in \eqref{eq:Z} according to whether the new label is absent or
present; the latter are $\{p\}\cup S$ with $P(S)\le t/p$, proving
\eqref{eq:Zrec}.  Since
$Z_V(y)=\1_{y\ge1}-F_V(y)$,
\[
 F_{V\cup\{p\}}(t)=F_V(t)+r\bigl(\1_{t/p\ge1}-F_V(t/p)\bigr).
\]
Integrate over $[X/4,X]$ and substitute $u=t/p$ in the last term.  The
indicator integral is exactly $H_X(p)$.  The final assertion follows from
Theorems \ref{thm:window} and \ref{thm:Eulerchar}.
\end{proof}

Equation \eqref{eq:Grec} is the sharp scalar obstruction.  If a child satisfies
$G_V(X/p)>H_X(p)$, adjoining the label contributes negatively.  Any proof that
replaces this correlated child by an unsigned reservoir loses the sign needed
at the root.

\section{The Mellin transform}

\begin{lemma}[Elementary logarithmic Mellin integral]\label{lem:logMellin}
For $n>0$ and $\Re s>0$,
\[
 \int_n^\infty \log(X/n)X^{-s-1}\,dX=\frac{n^{-s}}{s^2}.
\]
\end{lemma}
\begin{proof}
Put $X=ne^u$.  The integral becomes
$n^{-s}\int_0^\infty ue^{-su}\,du$.  Integration by parts gives $1/s^2$.
\end{proof}

Because $a(n)=O(1)$, one has $\A_X=O(\sqrt X)$; hence its Mellin transform
converges absolutely for $\Re s>1/2$.

\begin{theorem}[Exact zero-safe transform]\label{thm:Mellin}
For $\Re s>1/2$ and $z=s+1/2$,
\begin{equation}\label{eq:Mellin}
 \boxed{
 \int_1^\infty\A_X X^{-s-1}\,dX
 =\frac{1-4^{-s}}{s^2}
  \left[6-\frac{3(1-2^{-z})(2-2^{-z})}{\zeta(z)}\right].}
\end{equation}
\end{theorem}
\begin{proof}
Finite truncation followed by absolute convergence allows termwise use of
Lemma \ref{lem:logMellin}:
\[
 \int_1^\infty R_X X^{-s-1}dX=\frac{A(z)}{s^2}.
\]
The substitution $X=4Y$ shows that the transform of $R_{X/4}$ is
$4^{-s}A(z)/s^2$.  Insert the explicit $A(z)$ above.
\end{proof}

If $\rho$ is a zeta zero with $\Re\rho>1/2$, then at
$s_0=\rho-1/2$ none of the factors
\[
1-4^{-s_0},\qquad1-2^{-\rho},\qquad2-2^{-\rho}
\]
vanishes: their possible zeros have real parts $0,0,-1$, respectively.
Thus \eqref{eq:Mellin} has a nonremovable pole at $s_0$ of the same order as
the zero of $\zeta$.

